\chapter{Funciones Generatrices}

\section{Funciones generatrices}

\begin{definicion}
La \textbf{función generatriz} de la sucesión $a_0, a_1, \ldots$ es
\[
f(x) = \sum_{i=0}^{\infty} a_i x^i.
\]
\end{definicion}

\begin{ejemplo}
$(1+x)^n$ es la función generatriz de $\binom{n}{0}, \binom{n}{1}, \ldots, \binom{n}{n}, 0, 0, \ldots$
\end{ejemplo}

\begin{ejemplo}
$\dfrac{1}{1-x} = 1 + x + x^2 + x^3 + \cdots$ es la función generatriz de $1, 1, 1, \ldots$
\end{ejemplo}

\begin{ejemplo}
$\dfrac{1}{(1-x)^2} = 1 + 2x + 3x^2 + \cdots$ es la función generatriz de $1, 2, 3, \ldots$
\end{ejemplo}

\begin{ejemplo}
La función generatriz de $0^2, 1^2, 2^2, \ldots$ es $\dfrac{x(1+x)}{(1-x)^3}$.
\end{ejemplo}

\begin{table}[h]
\centering
\caption{Funciones generatrices fundamentales}
\begin{tabular}{|c|l|}
\hline
\textbf{No.} & \textbf{Función generatriz} \\
\hline
1 & $(1+x)^m = \sum_{n=0}^m \binom{m}{n} x^n$ \\
2 & $(1+ax)^m = \sum_{n=0}^m \binom{m}{n} a^n x^n$ \\
3 & $(1+x^m)^n = \sum_{k=0}^n \binom{n}{k} x^{mk}$ \\
4 & $\frac{1-x^{m+1}}{1-x} = 1+x+\cdots+x^m$ \\
5 & $\frac{1}{1-x} = \sum_{n=0}^{\infty} x^n$ \\
6 & $\frac{1}{(1-x)^m} = \sum_{n=0}^{\infty} \binom{m+n-1}{n} x^n$ \\
7 & $\frac{1}{(1+x)^m} = \sum_{n=0}^{\infty} (-1)^n \binom{m+n-1}{n} x^n$ \\
8 & $e^x = \sum_{n=0}^{\infty} \frac{x^n}{n!}$ \\
9 & $\ln(1+x) = \sum_{n=1}^{\infty} (-1)^{n+1}\frac{x^n}{n}$ \\
10 & $\frac{1}{(1-x)^2} = \sum_{n=0}^{\infty} (n+1)x^n$ \\
\hline
\end{tabular}
\end{table}

\begin{ejemplo}
Mónica compró 12 naranjas para Graciela, María y Francisco. Graciela recibe al menos 4, María al menos 2 y Francisco entre 2 y 5. La función generatriz es:
\[
f(x) = (x^4 + \cdots + x^8)(x^2 + \cdots + x^6)(x^2 + \cdots + x^5).
\]
El coeficiente de $x^{12}$ es $14$.
\end{ejemplo}